% Generated draft: requires mathematical and source review.
\documentclass[11pt]{article}
\usepackage[margin=25mm]{geometry}
\usepackage{fontspec}
\setmainfont{Latin Modern Roman}
\usepackage{amsmath,amssymb,mathtools,unicode-math}
\setmathfont{Latin Modern Math}
\usepackage{enumitem,needspace,xurl,hyperref}
\setlength{\parindent}{0pt}
\setlength{\parskip}{5pt}
\setlength{\emergencystretch}{3em}
\newcommand{\N}{\mathbb{N}}
\newcommand{\Z}{\mathbb{Z}}
\newcommand{\Q}{\mathbb{Q}}
\newcommand{\R}{\mathbb{R}}
\newcommand{\C}{\mathbb{C}}
\newcommand{\F}{\mathbb{F}}
\newcommand{\A}{\mathbb{A}}
\newcommand{\PP}{\mathbb{P}}
\newcommand{\E}{\mathbb{E}}
\newcommand{\eps}{\varepsilon}
\newcommand{\dd}{\,\mathrm{d}}
\begin{document}
\section*{UnsolvedMath ID: 800583}
\textbf{Source problem code:} catalogue entry
\subsection{Definitions and mathematical statement}
Let \(\mathfrak c=|\R|=2^{\aleph_0}\). For a set \(X\subseteq\R\), use its subspace topology. A countable Borel \(\omega\)-cover of \(X\) is a countable family \(\mathcal U\) of Borel subsets of \(X\) such that \(X\notin\mathcal U\) and every finite subset of \(X\) is contained in some member of \(\mathcal U\). A countable Borel \(\gamma\)-cover of \(X\) is a countably infinite family \(\mathcal U\) of Borel subsets of \(X\), with \(X\notin\mathcal U\), such that each point of \(X\) belongs to all but finitely many members of \(\mathcal U\). Write \(\mathrm B_\Omega(X)\) and \(\mathrm B_\Gamma(X)\) for these two classes of covers.

The selection property \(\mathsf S_1(\mathrm B_\Omega,\mathrm B_\Gamma)\) means that for every sequence \((\mathcal U_n)_{n\in\N}\) with \(\mathcal U_n\in\mathrm B_\Omega(X)\), there are choices \(U_n\in\mathcal U_n\) such that
\[
\{U_n:n\in\N\}\in\mathrm B_\Gamma(X).
\]
In particular, the selected family must be infinite; the displayed condition concerns the family of distinct selected sets.

For this rendering, Martin's Axiom \(\mathrm{MA}\) denotes the assertion that for every partially ordered set satisfying the countable chain condition and every family of fewer than \(\mathfrak c\) dense subsets, there is a filter meeting every member of that family. The countable chain condition means that every antichain is countable. With stronger conditions ordered below weaker ones, a filter is a nonempty upward closed, downward directed subset, and a set \(D\) is dense if every condition has a stronger condition in \(D\).

\textbf{Question.} Does Martin's Axiom imply, over \(\mathrm{ZFC}\), that there exists \(X\subseteq\R\) such that
\[
|X|=\mathfrak c
\qquad\text{and}\qquad
X\models\mathsf S_1(\mathrm B_\Omega,\mathrm B_\Gamma)?
\]
Here the last expression abbreviates the selection property defined above, rather than satisfaction in a first-order structure.

\textbf{Unresolved source qualifications.} The supplied excerpt does not define its cover notation or state whether its convention for Martin's Axiom includes \(\neg\mathrm{CH}\). The definitions above give an explicit conventional reading, with no additional assumption of \(\neg\mathrm{CH}\); agreement with the source conventions requires review. The excerpt precedes the question with the assertion that all constructions mentioned earlier can be carried out using Martin's Axiom, except perhaps ``Setting (n).'' Neither those constructions nor Setting (n) is supplied, so their relation to this question cannot be specified. The attribution marker \([\mathrm{MilNonGamma}]\) is retained but is not expanded by the supplied record.
\subsection{Short English statement}
Under the stated conventional definitions, does Martin's Axiom guarantee a set of reals of size continuum from whose every sequence of countable Borel \(\omega\)-covers one can select one member per cover so that the selected family is a \(\gamma\)-cover? The registry labels this problem open; that label is supplied source data, not an independently verified assessment.
\subsection{Sources}
\begin{itemize}
\item Local UnsolvedMath registry, permanent ID \(800583\), exact source problem code \texttt{catalogue entry}; upstream ID \(800583\), revision \texttt{d371ee917ac9647fbb06cd6fc529deec6f3dda82}. The record lists Tsaban as proposer, gives the year \(2006\), and labels the problem open.
\item Tsaban, \emph{Selection Principles and Special Sets of Reals: Open Problems} (2006), Problem 9. Supplied source URL: \url{https://arxiv.org/abs/math/0609601}. The record reports access on 2026-10-16 and extraction by \texttt{arxiv\_latex}; these are supplied provenance details, not a claim of independent source access or verification.
\item Supplied attribution/context marker: \texttt{[MilNonGamma]}. No corresponding bibliographic entry is included in the record. The source's contextual qualification about earlier constructions and the possible exception of ``Setting (n)'' is preserved above; the referenced material is absent.
\end{itemize}
\end{document}
